% TOP_PROBLEM: {"id": 30006821, "title": "Equivalence or Strict Containment of QMA and QMA(2)", "category_id": 16, "category": {"id": 16, "name": "physics", "display_name": "Mathematical Physics", "description": "Problems at the intersection of mathematics and physics.", "slug": "physics", "order_index": 16, "created_at": "2026-07-31T15:26:25.671Z"}, "status": "open"}
% FIRST_POSTED: 2026-09-25
% ENTRY_KIND: statement_only
% !TeX program = lualatex
% Definition and sources only; this file is not an LLM solution attempt.
\documentclass[11pt]{article}
\usepackage[margin=25mm]{geometry}
\usepackage{fontspec}
\setmainfont{Latin Modern Roman}
\usepackage{amsmath,amssymb,mathtools}
\usepackage{unicode-math}
\setmathfont{Latin Modern Math}
\usepackage{xurl,hyperref}
\hypersetup{colorlinks=true,urlcolor=blue}
\setcounter{secnumdepth}{0}
\setlength{\parindent}{0pt}
\setlength{\parskip}{5pt}
\setlength{\emergencystretch}{3em}
\begin{document}
\section*{\texorpdfstring{Equivalence or Strict Containment of QMA and QMA(2)}{Equivalence or Strict Containment of QMA and QMA(2)}}\label{30006821}
\subsection{Definitions and mathematical statement}
For a promise problem $(Y,N)$, a QMA verifier is a uniform polynomial-time quantum circuit given one polynomial-size quantum witness, with completeness at least $2/3$ and soundness at most $1/3$. QMA(2) instead receives two registers whose joint witness must be a product $\rho\otimes\sigma$. Each yes instance has an accepted product witness, while soundness is required for every product witness on a no instance. Determine whether
\[
 \mathrm{QMA}(2)=\mathrm{QMA}
\]
or the inclusion of QMA in QMA(2) is strict. No behavior is required outside the promise. The absence of entanglement is across the two witness registers, not within each register, and the model has no advice or oracle.
\par\smallskip\textit{Formulation and concept references:} \href{https://theoryofcomputing.org/articles/v005a001/v005a001.pdf}{[S1]}.

\subsection{Short English statement}
Does allowing two quantum witnesses promised to be unentangled give more verification power than allowing one unrestricted quantum witness?
\subsection{Sources}
\begin{itemize}
\item[S1]Scott Aaronson, Salman Beigi, Andrew Drucker, Bill Fefferman and Peter Shor, The Power of Unentanglement, Theory of Computing 5 (2009), 1-42, DOI 10.4086/toc.2009.v005a001.
\url{https://theoryofcomputing.org/articles/v005a001/v005a001.pdf}.
\textit{Formulation source.}
\end{itemize}
\end{document}
